\chapter{光合作用：植物不是把阳光直接变成食物}

\begin{kaichang}
今天早上的面包、米饭或者馒头，先别急着咽下去，想一想：它原本是一株小麦或水稻。再往前推，这株植物是从一粒几克重的种子长成的——一株成熟的小麦比种子重几百倍。

多出来的这几百倍质量，是从哪来的？

先猜，写在纸上。候选答案有四个：从土壤里“吃”来的；从水里来的；把阳光“变成”了身体；从空气里来的。我猜不少人会选土壤——植物不是扎根在土里吗？也可能有人选阳光——课本上不是说“植物靠阳光制造食物”吗？

这一章结束时，我们回来对答案。你会发现，正确答案既在四个选项之中，又和大多数人理解的方式完全不同：阳光确实是植物的命脉，但它不是“原料”，而是“工资”。植物没有、也无法把光直接变成自己的身体。
\end{kaichang}

\section{一株柳树称出的质量之谜}

先说一个真实做过的实验。十七世纪四十年代，比利时医生海尔蒙特（Jan Baptista van Helmont）找了一只大木桶，烘干土壤后称得 \ud{90}{kg}，种下一株 \ud{2.3}{kg} 的柳树苗，盖上盖子防止尘土落入。此后五年，他只浇雨水和蒸馏水。五年后，柳树长到 \ud{76}{kg} 多，而把土重新烘干一称：只少了约 \ud{60}{g}。

七十多千克的新木头和枝叶，土壤只贡献了六十克。海尔蒙特的结论是：植物的身体“几乎全部来自水”。这个结论错了一半——水确实是原料之一，但最大的那部分质量，他漏算了：空气。他那个时代还没有“气体也有质量、还能参与化学反应”的概念，这不能怪他。但这个实验留下的定量结论至今成立：\textbf{土壤不是植物身体的主要来源}。

那阳光呢？光是能量，不是物质。你没办法把光“称重”后缝进木头里——能量没有质量账本上这一栏（至少在日常的化学尺度上，光合作用涉及的质量变化小到完全测不出来）。所以如果植物的身体不是土做的、也不是光变的，那它只能来自剩下的地方：水和空气。后面我们会看到，一片叶子干重的九成以上，追根溯源是二氧化碳和水。你呼吸吐出的“废气”，正是植物盖房子的砖头。

还有一个你在池塘边就能看见的现象：晴天的中午，水草叶片上挂满了小气泡，一串串往上冒。气泡里是什么？它从哪来？为什么阴天就几乎不冒？把这几个问题记住，它们是本章的路标。

\section{镜头推进：叶子里的绿色工厂}

把镜头推进到叶片内部。叶片细胞里有一种绿色的小颗粒，叫\textbf{叶绿体}，直径只有几微米，一个叶肉细胞里装着几十个。第 51 章讲过细胞膜是磷脂双层围出的边界，叶绿体更进一步：它有两层膜，内部还叠着一摞摞扁平的“硬币”，叫类囊体，硬币膜上密密麻麻嵌着两类东西——吸光的色素分子和一系列蛋白质机器。

绿色从哪来？类囊体膜上的主角是叶绿素——第 10 章提过，它的分子正中央卡着一个镁原子，周围一圈碳氮环，拖着一条长长的疏水尾巴，这条尾巴把它锚在膜上。叶绿素对光非常“挑食”：它强烈吸收蓝紫光（约 $430$ 纳米附近）和红光（约 $660$ 纳米附近），而对绿光（约 $500$～$600$ 纳米）几乎不理睬。绿光大部分被反射和透射出来，所以叶子在你眼里是绿的。\textbf{叶子的颜色，恰恰写着它“不吃”哪种光}。

为什么是特定的波长？回到第 9 章的规则：光是一份一份的，每份（一个光子）的能量由波长决定，波长越短，单份能量越高；而分子里的电子只接受“恰好够格”的能量台阶。蓝紫光和红光光子的能量，正好对准叶绿素分子里某个电子的跃迁台阶——光子被吸收，电子被“踢”上高能级。绿光的能量对不上号，就被放行了。这就像投币售货机只认特定面额的硬币，不是它“喜欢”，是结构决定了匹配。

一个被踢上高能级的电子是不安稳的。在许多物质里，这样的电子很快就会跌回原位，把能量以热或微弱荧光的形式还回去——防晒霜里的分子就是这么干掉的紫外线。但叶绿体干了一件完全不同的事：\textbf{趁电子还没跌回去，把它“截走”}，送上一条电子传递链。这才是光合作用的真正起点——不是“吸收阳光”，而是“用阳光把电子从低能级撬到高能级，再收进分子的口袋”。

\begin{shiyan}
\textbf{观察：水草在光下冒气泡}

\textbf{材料}：一枝金鱼藻（花鸟市场或水族店有售，也可用新鲜的绿萝枝条代替，气泡会少些）、透明玻璃杯、清水、一盏台灯、一把剪刀、一个计时器。

\textbf{步骤}：①在水中剪下一段约 10 厘米的水草（在水中剪，避免切口进入空气）；②切口朝上放入杯中，加水没过；③把台灯放在距杯子约 20 厘米处照光，计时观察切口处冒出的气泡，数一数 5 分钟冒多少个；④关灯或把杯子移到暗处，同样数 5 分钟。

\textbf{观察点}：光照下气泡明显增多，黑暗里几乎停止；气泡都从那一个切口冒出，而不是整片叶子随便冒——说明气体是植物体内“生产”出来的，不是水里溶解的空气受热跑出来的。

\textbf{安全提示}：实验本身没有危险品，但台灯开久了灯罩发烫，别用手摸；水杯和台灯保持距离，小心水泼溅到电器。想知道气泡是什么气体，需要用化学方法检验——那一步请在学校实验室里由老师带着做，不要自己收集气体点火测试。
\end{shiyan}

\section{一场接力赛：从水分子到 ATP}

现在把镜头推到最近，看这团被截住的高能电子怎么干活。整个过程像一场接力赛，全程发生在类囊体膜上和膜附近，可以分成两段：前一段直接需要光，叫\textbf{光反应}；后一段不直接需要光，叫\textbf{碳反应}（教科书常叫“暗反应”，但这个名字有误导，它白天黑夜都在跑，我们不用它）。

第一棒：水分子被拆开。叶绿素失去电子后变成了“缺电子”的状态，必须立刻补回一个，否则机器就停转。补给站是水。类囊体膜上有一组含锰的蛋白质机器，把水分子的电子硬生生拽下来补给叶绿素。水被拆开后，氢变成氢离子 \chem{H^+} 留在膜内，而两个水分子里的两个氧原子两两配对，变成氧气分子 \chem{O_2} 释放出去——你刚才在水草切口看到的气泡，源头就在这里。

请在这一句上停三秒钟，因为它是本章的第一条硬知识：\textbf{植物释放的氧气，来自水，不是来自二氧化碳}。水才是被拆开的那个，二氧化碳从头到尾没有贡献一个氧原子给空气。这不是猜测，是有同位素实验铁证的结论，证据故事稍后讲。

第二棒：电子下楼梯，质子被抽上楼。从水里来的高能电子并不直接用于“存起来”，而是顺着膜上的一系列蛋白质接力棒逐级传递，每传一级，能量降一点——和第 53 章讲的线粒体电子传递链几乎是同一个剧本。降下来的能量被用来把 \chem{H^+} 从膜的“外面”抽到“里面”（类囊体腔），在膜两侧攒出一个又浓又酸的质子梯度。然后，质子顺着梯度回流，推动膜上一台纳米级旋转马达——ATP 合酶——把 ADP 拧成 ATP。化学渗透这套机制你在第 53 章已经见过一次，这里是它的第二次登场：进化把同一招用在了两套系统里，一个是“拆糖放能”，一个是“借光充能”。

第三棒：把电子存进“活期存折”。电子传到链条末端，和 \chem{H^+} 一起交给一种叫 \chem{NADP^+} 的电子载体，变成 NADPH。NADPH 是第 50 章讲过的 NADH 的“近亲”，同样是随身携带高能电子的运钞车。

光反应到此结束。清点战利品：氧气（副产品，放出去了）、ATP（能量现金）、NADPH（高能电子存折）。注意，到现在为止，\textbf{一个碳原子都还没固定}。光的作用已经全部完成——它换来的不是糖，而是两种“高能货币”。

\section{碳反应：用现金和存折买碳}

接下来的事发生在叶绿体的液态基质里，不再直接需要光：植物要用刚才挣来的 ATP 和 NADPH，把空气中的二氧化碳“焊”进有机物。这条焊接流水线是一圈循环，叫\textbf{卡尔文循环}，大致三步走。

第一步，抓碳。空气中 \chem{CO_2} 的浓度只有约万分之四，稀薄得很。叶绿体里有一种酶，负责抓住一个 \chem{CO_2}，把它接到一个五碳糖分子上，再立刻裂成两个三碳分子。这个酶叫 Rubisco——名字拗口，但值得记住，因为它是地球上数量最多的蛋白质：全世界的 Rubisco 加在一起，总质量以亿吨计。为什么这么舍得下本？因为它干活慢，一秒只能抓几个 \chem{CO_2}（一般的酶一秒能处理成百上千个底物），植物只能靠堆数量弥补速度。你看到的每一片绿色，里面有相当大一部分其实是 Rubisco 的颜色。

第二步，还原。三碳分子先后花掉 ATP 的化学能和 NADPH 的高能电子，被加工成一种三碳糖。光反应挣来的现金和存折，就在这一步花光——这也正是两段反应的连接点：光反应供币，碳反应消费。

第三步，续跑。每抓 3 个 \chem{CO_2} 循环若干圈，净赚 1 个三碳糖离开循环，其余的三碳分子再花 ATP 重新拼回五碳糖，继续抓碳。离开循环的三碳糖汇成葡萄糖，再按第 46 章讲的方式聚成淀粉存起来，或者运去建造纤维素、蛋白质和脂肪。

回头看本章标题那句话：植物不是把阳光直接变成食物。真实的因果链是——光能先把水拆开放出氧气，顺手把电子撬上高能级；高能电子的能量被分批兑换成 ATP 和 NADPH；这两种货币再驱动酶把 \chem{CO_2} 焊接、还原成糖。\textbf{糖的身体来自二氧化碳和水，糖里的能量才是阳光的遗产}。光是发电厂的电，不是盖楼的砖。

\section{符号登场：一个必须写“十二水”的反应式}

忍到现在，终于可以写反应式了。许多教材把光合作用总反应简写成：

\[\chem{6CO_2} + \chem{6H_2O} \rightarrow \chem{C_6H_{12}O_6} + \chem{6O_2}\]

这个写法没有错，但它藏起了一个关键事实。数一氧原子：左边 \chem{CO_2} 带了 12 个氧，水里 6 个氧，共 18 个；右边葡萄糖占掉 6 个，\chem{O_2} 占 12 个。看起来“6 个水给 6 个 \chem{O_2} 供氧”也说得通？这正是简写版容易让人误会的陷阱。既然氧气全部来自水，水就必须在左边出现\textbf{两次}——一次作为被拆开的电子供体，一次作为反应环境。更诚实的写法是：

\[\chem{6CO_2} + \chem{12H_2O} \rightarrow \chem{C_6H_{12}O_6} + \chem{6O_2} + \chem{6H_2O}\]

左边 12 个水分子里，有 12 个氧原子全部进了右边的 \chem{6O_2}；葡萄糖和右边新出现的 6 个水分子里的氧，则全部来自 \chem{CO_2}。而光反应里真正被拆的那一步，符号长这样：

\[\chem{2H_2O} \rightarrow \chem{O_2} + \chem{4H^+} + \chem{4e^-}\]

请记住这个不对称：每放出 1 个 \chem{O_2}，要拆 2 个水、拽走 4 个电子。这 4 个电子就是整条流水线的“原始资本”。

\begin{zhengju}
凭什么说氧气来自水而不是二氧化碳？两边的原子都一样是氧，普通实验根本分不出来——这就像两杯一模一样的水混在一起，你怎么证明某个水分子来自哪一杯？

给分子“贴标签”的办法，是同位素。第 7 章讲过，同一元素可以有质子数相同、中子数不同的版本。普通的氧是氧-16，还有一种稳定的氧-18，重 2 个中子，化学行为几乎一样，但可以用仪器区分。二十世纪三十年代，美国斯坦福大学的鲁本（Samuel Ruben）和卡门（Martin Kamen）团队想到了主意：分别制备“普通水 + 氧-18 标记的 \chem{CO_2}”和“氧-18 标记的水 + 普通 \chem{CO_2}”两套培养液，让绿藻在里面进行光合作用，然后检测放出来的 \chem{O_2} 带不带标记。

结果干净利落：\chem{CO_2} 被标记时，放出的氧气是普通氧；水被标记时，放出的氧气带上了氧-18。氧气里氧原子的“户籍”，百分之百登记在水名下。1941 年这组实验发表，成为光合作用的基石之一。

其实早在实验之前，已经有人从另一条路猜到了答案。荷兰裔微生物学家范尼尔（C. B. van Niel）研究紫硫细菌——一种能进行光合作用但\textbf{不放氧}的细菌，它们分解硫化氢 \chem{H_2S} 而不是水，产物是硫磺颗粒而不是氧气。范尼尔注意到：细菌分解 \chem{H_2S} 放硫，和植物分解 \chem{H_2O} 放氧，剧本一模一样，只是“被拆的分子”和“拆下来的残渣”换了。于是他推断：植物的氧也该来自被拆的那个分子——水。一个从边缘生物来的类比，十年后被同位素实验证实。科学里很多大结论，都是“猜想先走，证据后到”，但能不能站住，永远由证据说了算。
\end{zhengju}

\section{光合作用和呼吸作用：像，但不是互为逆过程}

现在处理一个被无数简写反应式喂大的误解。把光合作用总式和第 53 章的细胞呼吸总式放在一起：

\begin{itemize}[itemsep=2pt]
\item 光合作用：\chem{6CO_2} + \chem{12H_2O} $\rightarrow$ \chem{C_6H_{12}O_6} + \chem{6O_2} + \chem{6H_2O}
\item 细胞呼吸：\chem{C_6H_{12}O_6} + \chem{6O_2} $\rightarrow$ \chem{6CO_2} + \chem{6H_2O}
\end{itemize}

太像了，对吧？箭头一反，两式几乎互为镜像。于是“光合作用是呼吸作用的逆反应”成了流传极广的说法。

\begin{wuqu}
“光合作用就是呼吸作用倒过来跑，植物白天光合作用、晚上呼吸作用。”

两个半句都错。第一，\textbf{路径完全不同}。逆着写一个总反应式很容易，但细胞里的化学反应是一步一步走的，走错一步就是另一种物质。呼吸作用在细胞质和线粒体里，用糖酵解、柠檬酸循环、电子传递链这一条路线（第 52、53 章）；光合作用在叶绿体里，用光系统、卡尔文循环这条完全不同的路线，两套装配线上的酶几乎不重叠。打个比方说：从北京到上海，高铁和飞机的起终点相同，但你不能说“坐飞机就是倒着坐高铁”。连原子的走法都不一样——呼吸释放的 \chem{CO_2} 里，氧原子来自葡萄糖和水；光合释放的 \chem{O_2} 里，氧原子来自水。同位素标记把每一方的原子去向查得清清楚楚，根本不是同一批原子在来回倒腾。

第二，\textbf{植物白天黑夜都在呼吸}。叶子造糖，但叶子的细胞也要吃饭——维持膜、合成蛋白质、修复损伤，样样要花 ATP，这些 ATP 大部分还是靠线粒体烧糖来的。白天你测到植物“放氧气”，只是因为光合放氧的速度超过了呼吸耗氧的速度，是两条流水线\textbf{同时开工、互相轧差}之后的净结果。晚上光合停机，呼吸继续，植物就和动物一样吸氧放二氧化碳。卧室里那盆绿萝夜里“抢氧气”？别担心，一盆植物的呼吸功率远不如和你一起睡觉的一个人——甚至不如一只猫。
\end{wuqu}

这两条流水线真正深层的联系，不在反应式的对称，而在\textbf{分子机器的对称}：两边都有电子传递链，都用质子梯度攒势，都用 ATP 合酶收账。第 53 章的化学渗透和本章的光反应，像同一套发动机装在两辆车上。这不是巧合，而是进化史上的一次“技术复用”——光合细菌和呼吸细菌的共同祖先，很早就发明了用膜和质子梯度存取能量的办法。这段渊源要到讲进化的第十篇再细说，这里先埋在心里。

\section{同一个配方，三种活法}

卡尔文循环是几乎所有绿色植物通用的“焊接车间”，但车间开工的环境千差万别。核心麻烦出在 Rubisco 身上：这位劳模有个职业病——它分不清 \chem{CO_2} 和 \chem{O_2}。本该抓二氧化碳的活性位点，时不时会抓一个氧气分子进去，产生的半成品不但没用，还得花 ATP 去“返工回收”，这个过程叫光呼吸。气温越高、叶片为了保水把气孔关得越紧，叶内 \chem{CO_2} 越稀、\chem{O_2} 越浓，抓错的概率就越大。炎热干旱的夏天中午，水稻小麦这类植物可能有相当一部分固定好的碳被光呼吸吐掉，等于工厂一边生产一边返修。

面对同一道难题，植物在进化中交出了三份答卷：

\begin{itemize}[itemsep=2pt]
\item \textbf{C$_3$ 植物}（水稻、小麦、大豆、绝大多数树木）：直接用卡尔文循环，固定出的第一个稳定产物是三碳分子，所以叫 C$_3$。在温凉湿润的环境里它们最经济——不搞任何额外建设；但高温干旱下被光呼吸拖累。
\item \textbf{C$_4$ 植物}（玉米、甘蔗、高粱）：给车间加了一道“前置泵”。先在叶肉细胞里用一种不挑食的酶把 \chem{CO_2} 抓进一个四碳分子（所以叫 C$_4$），再把四碳分子运到叶片深处的维管束鞘细胞里“卸货”——那里的 \chem{CO_2} 被浓缩了好几倍，Rubisco 在二氧化碳的包围中工作，几乎不会抓错。代价是每固定一个碳要多花 ATP，相当于给浓缩泵交电费。在强光和炎热下，这笔电费花得值。
\item \textbf{CAM 植物}（仙人掌、菠萝、龙舌兰）：换一个思路——不在空间上隔离，在时间上错开。白天气孔紧闭，一粒水也不放出去；夜里凉爽时气孔大开，吸进的 \chem{CO_2} 先存成苹果酸，天亮后关起门来慢慢拆解，供卡尔文循环使用。节水做到极致，代价是生长缓慢——沙漠里的植物不赶时间。
\end{itemize}

三份答卷没有高下之分，只有“和谁的环境匹配”。这是自然选择的典型作风：不追求完美设计，只在现有零件上修修补补，够用就行。Rubisco 的“职业病”至今没被修好——不是不想修，是修它的每一步中间状态都比现在更差，进化翻不过那道山谷。

\section{这个模型什么时候会失灵}

本章的图景很好用，但请记住三条边界。

第一，\textbf{“光合作用”不等于“放氧气”}。范尼尔的紫硫细菌就是反例：它们有光合色素、有电子传递、能固定碳，但拆的是 \chem{H_2S} 而不是水，不放氧。放氧光合作用（拆水）只是光合作用大家族里的一支，只是恰好这一支改变了整个地球——今天大气里五分之一的氧气，全是它几十亿年拆水的“工业废气”。第 56 章讲早期生命时，这场“大氧化事件”还会回来。

第二，\textbf{总反应式只写净结果，隐藏了速率}。“光够、水够、二氧化碳够，糖就长出来”——这是平衡账，不是速度账。真实叶片的生产速度受限于最慢的那一环：弱光时卡在光反应，强光下卡在 Rubisco 的慢手慢脚，高温时卡在气孔关闭。温室种植正是顺着这个逻辑操作：给温室补充 \chem{CO_2}（把浓度从万分之四提到万分之八左右）、延长补光，让“最慢一环”换人，产量就能明显上升——这是把本章模型直接变成农业技术的例子。

第三，\textbf{效率远没有想象中高}。晒在叶片上的太阳能，最终存进糖里的通常只有百分之一到百分之几。其余的去向：绿光和红外光直接没被吸收；吸收了的能量一部分变成热散掉；光呼吸再吃掉一块。觉得植物“效率低”？先别急着嫌弃——它用的是免费的光、空气和水，零燃料成本运行了二十多亿年，还顺手养活了整个生物圈。人类的光伏板效率虽高，工厂和电费的账本可一点都不便宜。

\begin{qiaoliang}
来算一笔账：一平方米的旺盛叶片，一天能造多少糖、放多少氧？实测的典型净光合速率约为每平方米每秒固定 $15$ 微摩尔（$15\times 10^{-6}$ mol）\chem{CO_2}。一个夏日按 $10$ 小时有效强光算：

\[
15\times 10^{-6} \times 3600 \times 10 \approx 0.54\ \text{mol CO}_2/\text{天}
\]

由总反应式，每 6 个 \chem{CO_2} 换 1 个葡萄糖、放 6 个 \chem{O_2}，即每天约造 $0.09$ mol 葡萄糖——约 \ud{16}{g}，同时放出 $0.54$ mol 氧气，按常温下每摩尔约 $24$ 升算，约 $13$ 升。

现在跟你比一比：一个安静的中学生一天大约消耗 $500$ 多升氧气。也就是说，\textbf{你一天的呼吸，需要大约 40 平方米旺盛生长的叶片满负荷晒一天才能抵消}。一片足球场大小的草地，大约能平衡二十个人的呼吸——这还没算草地自己在夜里的呼吸消耗。下次看到“城市绿肺”这个词，你就知道它背后是一本实打实的原子账本。
\end{qiaoliang}

\section{回到那块面包}

现在对开场的那道题收账。一株小麦比种子重出来的几百倍质量，来自哪里？

不是土壤——海尔蒙特的桶已经替你排除了，土壤只贡献了矿物质零头。不是阳光直接变的——光没有变成任何一个原子，它被叶绿素收下，换成 ATP 和 NADPH，花掉，散场。真正的答案是：\textbf{绝大部分来自空气里的二氧化碳，配上水}。碳是砖，水是泥，光是电。你吃下的那块面包，本质上是一块“凝固的空气”加上一截“捕获的阳光”——这句话写进作文可能被老师扣分，但按原子账本，它分毫不差。

水草切口冒出的气泡也有了解释：那是水被拆解时放出的氧，沿着植物体内的通气组织从伤口逸出。台灯是工资发放机，灯一关，拆水停工，气泡断流。

顺便回答一个你可能一直想问的问题：既然植物能把 \chem{CO_2} 变成糖，我们为什么不照抄一套人工装置来解决粮食和气候问题？科学家确实在试——这叫人工光合作用。难点不在原理，原理本章已经讲完；难点在工程：水是最难拆的稳定分子之一，要廉价、耐用、高效地拆它，人类目前的催化剂还远不如叶绿体里那台含锰机器“皮实”。大自然用二十多亿年迭代出的方案，没那么容易抄。

\section{糖造出来之后呢}

这一章我们只跟到了葡萄糖出厂。但一个活体细胞不可能只靠囤淀粉过日子：糖要流去烧掉换 ATP（第 52、53 章），要变成纤维素砌墙，要转成脂肪存进种子，要为蛋白质和核酸提供碳骨架。绿色工厂的车间之间不是单行道，而是一张处处相通、互相调节的路网——胰岛素怎样指挥全身细胞分配燃料，饥饿时身体改烧什么，都在这张网里。下一章，我们就把镜头从单条流水线拉远，看这张代谢网络的全貌。

\begin{tiaozhan}
固定碳的账到底多贵？卡尔文循环每净赚 1 个三碳糖（3 个碳），要固定 3 个 \chem{CO_2}，花掉 9 个 ATP 和 6 个 NADPH。这些“货币”折合多少个光子？电子的走法是：每从水拽出 4 个电子（放出 1 个 \chem{O_2}），这些电子要先后被两个光系统各“踢”一次，也就是至少 $8$ 个光子，产生的电子流大约换回 3 个 ATP 和 2 个 NADPH。按比例配平，固定 3 个 \chem{CO_2} 需要的货币大致对应 $30$ 个上下的光子。再算能量：30 个红光光子的总能量约 $5100$ 千焦每摩尔，而 3 个碳存成糖能保存的化学能约 $1400$ 千焦——光反应加碳反应的“账面效率”大约三成。看起来浪费？别忘了这几乎是一切化学生物过程的天花板：能量每转一次手，就有一份变成热散掉，这是热力学第二定律收的过路费，谁也逃不掉。跳过本框不影响主线。
\end{tiaozhan}

\section*{练一练}

\begin{sikaoti}
\item 画一张物质流图：用方框和箭头画出“\chem{CO_2}、水、光能、氧气、ATP、NADPH、葡萄糖”在光合作用中的来龙去脉，标出每一步发生在叶绿体的哪个位置（类囊体膜上还是基质里）。图旁注明：方框和箭头是人为的表示法。
\item 一位同学用“普通水 + 氧-18 标记的 \chem{CO_2}”培养绿藻，检测发现葡萄糖和新生成的水里有氧-18，而放出的 \chem{O_2} 里没有。这个实验结果支持“氧气来自水”还是反驳它？请说明推理。
\item 把一个绿色植物密封在透明玻璃罩里，光照足够，几天之后罩内的 \chem{CO_2} 浓度会降低、\chem{O_2} 浓度会升高，但最终会稳定在某个水平不再变化。用“光合与呼吸同时开工、互相轧差”的观点解释为什么会稳定，并预测：如果把罩子移到完全黑暗处，气体会怎么变？
\item 玉米（C$_4$ 植物）和水稻（C$_3$ 植物）种在同一块炎热干旱、阳光强烈的田里，哪一种更可能高产？为什么？如果把它们移到阴凉湿润的环境，你的预测会变吗？
\item 本章说“光是电，不是砖”。请用这句话的思路反驳一个说法：“多晒太阳，植物就长得壮，所以植物是吃阳光长大的。”你的反驳里要出现至少三个本章讲过的分子或结构名称。
\item 估算题：你家阳台有一盆约 0.1 平方米叶面积的绿萝。参照本章“原理桥梁”里的账本，估算它旺盛生长一天大约放多少升氧气，并说明为什么这点氧气不值得你为它“抢空气”而担心。
\end{sikaoti}
